% Auto-generated by repro/analysis/numbers.py -- do not edit by hand.
% results: repro/results (runs.csv sha256 49d2b357b63396a6); QALM config a; threshold 32%

% --- setup
\newcommand{\numSuiteCircuits}{248} % circuits in repro/inputs/S248.csv (every S248 column must have them all) [paper: 248 at texts/0-abstract.tex:38; texts/1-introduction.tex:121; texts/5-eval.tex:27,40]
\newcommand{\numSuiteSourceCircuits}{250} % *.qasm files in the GUOQ benchmark directory S248 is taken from [paper: 250 at texts/5-eval.tex:28]
\newcommand{\numSuiteRemovedCircuits}{2} % source circuits not in S248 (['qft_10', 'qft_16']) [paper: 2 at texts/5-eval.tex:28-29 (qft_10, qft_16)]
\newcommand{\numGateCountOrders}{five} % floor(log10(max/min original total)) = floor(log10(768096/5)), as a word [paper: five at texts/5-eval.tex:40]
\newcommand{\numJobs}{3676} % runs.csv rows with a result (unique job ids)
\newcommand{\numHosts}{1} % distinct hosts in runs.csv (['qalm'])
\newcommand{\numPeakRssGiB}{31.9} % max over runs of the peak resident set size (GiB)
\newcommand{\numPeakCgroupGiB}{32.0} % max over runs of the cgroup memory.peak (GiB; includes page cache)
\newcommand{\numJobsRerun}{3} % runs whose reported attempt is not the first (memory-cap escalation, preemption or infrastructure failure)
\newcommand{\numMemCapGiB}{32} % largest memory cap tier (GiB): every reported run ran with headroom below it (smaller tiers escalate) [paper: 32 at texts/5-eval.tex:86]
\newcommand{\numJobCores}{124} % job cores 2-63,66-127 (one job per physical core)
\newcommand{\numMaxConcurrent}{124} % max jobs running at once on one host
\newcommand{\numMeanConcurrent}{96} % time-weighted mean of jobs running on one host while jobs were still being launched (first to last launch), averaged over hosts; rounded
\newcommand{\numMemEscalations}{0} % re-runs with the next memory-cap tier (requeue reason esc)
\newcommand{\numPreemptions}{0} % jobs preempted by the memory watchdog and re-run (reason pre)
\newcommand{\numInfraReruns}{3} % re-runs after infrastructure failures (reason infra)
\newcommand{\numCampaignHours}{42} % wall-clock hours from the first launch to the last event on a host

% --- S248 averages (tab:avg-comparison)
\newcommand{\numFidQalm}{0.4017} % mean fidelity 0.999^(total-cx)*0.99^cx over the circuits of column qalm (metrics.averages) [paper: 0.3864 at figtex/averages_table.tex:14-15; texts/5-eval.tex:124]
\newcommand{\numFidQalmOneMin}{0.3898} % mean fidelity 0.999^(total-cx)*0.99^cx over the circuits of column qalm_1min (metrics.averages) [paper: 0.3747 at figtex/averages_table.tex:14-15]
\newcommand{\numFidGuoqTwoQ}{0.3476} % mean fidelity 0.999^(total-cx)*0.99^cx over the circuits of column guoq_2q (metrics.averages) [paper: 0.3475 at figtex/averages_table.tex:14-15; texts/5-eval.tex:126]
\newcommand{\numFidGuoqTotal}{0.3692} % mean fidelity 0.999^(total-cx)*0.99^cx over the circuits of column guoq_total (metrics.averages) [paper: 0.3676 at figtex/averages_table.tex:14-15]
\newcommand{\numFidQuesoTwoQ}{0.3348} % mean fidelity 0.999^(total-cx)*0.99^cx over the circuits of column queso_2q (metrics.averages) [paper: 0.3362 at figtex/averages_table.tex:14-15]
\newcommand{\numFidQuesoTotal}{0.3544} % mean fidelity 0.999^(total-cx)*0.99^cx over the circuits of column queso_total (metrics.averages) [paper: 0.3554 at figtex/averages_table.tex:14-15]
\newcommand{\numFidVoqc}{0.3373} % mean fidelity 0.999^(total-cx)*0.99^cx over the circuits of column voqc (metrics.averages) [paper: 0.3428 at figtex/averages_table.tex:14-15; texts/5-eval.tex:125]
\newcommand{\numFidQuartz}{0.3514} % mean fidelity 0.999^(total-cx)*0.99^cx over the circuits of column quartz (metrics.averages) [paper: 0.3526 at figtex/averages_table.tex:14-15]
\newcommand{\numCxRedQalm}{25.59} % mean 2-qubit (CX) gate reduction % (circuits with cx>0) over the circuits of column qalm (metrics.averages) [paper: 24.61 at figtex/averages_table.tex:16-17; texts/5-eval.tex:116,123]
\newcommand{\numCxRedQalmOneMin}{21.71} % mean 2-qubit (CX) gate reduction % (circuits with cx>0) over the circuits of column qalm_1min (metrics.averages) [paper: 21.18 at figtex/averages_table.tex:16-17]
\newcommand{\numCxRedGuoqTwoQ}{24.78} % mean 2-qubit (CX) gate reduction % (circuits with cx>0) over the circuits of column guoq_2q (metrics.averages) [paper: 25.41 at figtex/averages_table.tex:16-17; texts/5-eval.tex:117]
\newcommand{\numCxRedGuoqTotal}{20.06} % mean 2-qubit (CX) gate reduction % (circuits with cx>0) over the circuits of column guoq_total (metrics.averages) [paper: 18.83 at figtex/averages_table.tex:16-17]
\newcommand{\numCxRedQuesoTwoQ}{12.93} % mean 2-qubit (CX) gate reduction % (circuits with cx>0) over the circuits of column queso_2q (metrics.averages) [paper: 13.19 at figtex/averages_table.tex:16-17]
\newcommand{\numCxRedQuesoTotal}{12.95} % mean 2-qubit (CX) gate reduction % (circuits with cx>0) over the circuits of column queso_total (metrics.averages) [paper: 13.37 at figtex/averages_table.tex:16-17]
\newcommand{\numCxRedVoqc}{4.06} % mean 2-qubit (CX) gate reduction % (circuits with cx>0) over the circuits of column voqc (metrics.averages) [paper: 4.24 at figtex/averages_table.tex:16-17; texts/5-eval.tex:124]
\newcommand{\numCxRedQuartz}{12.04} % mean 2-qubit (CX) gate reduction % (circuits with cx>0) over the circuits of column quartz (metrics.averages) [paper: 12.42 at figtex/averages_table.tex:16-17]
\newcommand{\numTotalRedQalm}{56.20} % mean total gate reduction % over the circuits of column qalm (metrics.averages) [paper: 52.34 at figtex/averages_table.tex:18-19; texts/5-eval.tex:117,122]
\newcommand{\numTotalRedQalmOneMin}{54.91} % mean total gate reduction % over the circuits of column qalm_1min (metrics.averages) [paper: 51.35 at figtex/averages_table.tex:18-19]
\newcommand{\numTotalRedGuoqTwoQ}{\ensuremath{-}1.97} % mean total gate reduction % over the circuits of column guoq_2q (metrics.averages) [paper: -2.63 at figtex/averages_table.tex:18-19; texts/5-eval.tex:118]
\newcommand{\numTotalRedGuoqTotal}{44.40} % mean total gate reduction % over the circuits of column guoq_total (metrics.averages) [paper: 43.89 at figtex/averages_table.tex:18-19]
\newcommand{\numTotalRedQuesoTwoQ}{27.34} % mean total gate reduction % over the circuits of column queso_2q (metrics.averages) [paper: 27.40 at figtex/averages_table.tex:18-19]
\newcommand{\numTotalRedQuesoTotal}{38.83} % mean total gate reduction % over the circuits of column queso_total (metrics.averages) [paper: 39.72 at figtex/averages_table.tex:18-19]
\newcommand{\numTotalRedVoqc}{45.05} % mean total gate reduction % over the circuits of column voqc (metrics.averages) [paper: 46.37 at figtex/averages_table.tex:18-19; texts/5-eval.tex:122]
\newcommand{\numTotalRedQuartz}{39.20} % mean total gate reduction % over the circuits of column quartz (metrics.averages) [paper: 40.73 at figtex/averages_table.tex:18-19]
\newcommand{\numFailQalm}{0} % runs of column qalm whose status is not OK ({})
\newcommand{\numFailQalmOneMin}{0} % runs of column qalm_1min whose status is not OK ({})
\newcommand{\numFailGuoqTwoQ}{3} % runs of column guoq_2q whose status is not OK ({'JVM_OOM': 3})
\newcommand{\numFailGuoqTotal}{4} % runs of column guoq_total whose status is not OK ({'JVM_OOM': 4})
\newcommand{\numFailQuesoTwoQ}{2} % runs of column queso_2q whose status is not OK ({'JVM_OOM': 2})
\newcommand{\numFailQuesoTotal}{4} % runs of column queso_total whose status is not OK ({'JVM_OOM': 4})
\newcommand{\numFailVoqc}{7} % runs of column voqc whose status is not OK ({'TIMEOUT': 7})
\newcommand{\numFailQuartz}{2} % runs of column quartz whose status is not OK ({'OOM_CAP': 2})
\newcommand{\numTotalRedMarginOneHour}{11.15} % total_red(qalm) - max over the six 1-h baselines of total_red (unrounded; the baseline is \numStrongestTotalRedBaseline) [paper: 5.97 at texts/1-introduction.tex:123]
\newcommand{\numTotalRedMarginOneMin}{9.86} % total_red(qalm_1min) - max over the six 1-h baselines of total_red (unrounded; the baseline is \numStrongestTotalRedBaseline) [paper: 4.98 at texts/1-introduction.tex:125]
\newcommand{\numStrongestTotalRedBaseline}{\voqc{}} % baseline column with the highest mean total gate reduction [paper: VOQC at texts/1-introduction.tex:124-125 ('the strongest baseline')]

% --- S248 one-minute QALM
\newcommand{\numOneMinFidDropRel}{2.9} % 100*(F(1 h) - F(1 min)) / F(1 h) (QALM, unrounded table means) [paper: <1 at texts/5-eval.tex:139-140 ('by less than 1\%')]
\newcommand{\numOneMinFidDropAbs}{0.0118} % F(1 h) - F(1 min) (QALM, unrounded table means) [paper: <1 at texts/5-eval.tex:139-140 ('by less than 1\%')]
\newcommand{\numOneMinCxRedDrop}{3.88} % cx_red(1 h) - cx_red(1 min), points (QALM, unrounded table means) [paper: <1 at texts/5-eval.tex:139-140 ('by less than 1\%')]
\newcommand{\numOneMinTotalRedDrop}{1.29} % total_red(1 h) - total_red(1 min), points (QALM, unrounded table means) [paper: <1 at texts/5-eval.tex:139-140 ('by less than 1\%')]
\newcommand{\numBudgetFactor}{60} % one-hour budget / one-minute budget (3600 s / 60 s) [paper: 60 at texts/5-eval.tex:139]

% --- S248 per-circuit win rates (fig:win_rate)
\newcommand{\numWinFidGuoqTwoQ}{91.5} % % of circuits where QALM >= guoq_2q on fidelity (metrics.win_rate win_or_tie; fidelity ties |d ln F| <= 1e-12) [paper: 83.9 at figure/win_rate.png; texts/0-abstract.tex:40; texts/1-introduction.tex:136]
\newcommand{\numStrictWinFidGuoqTwoQ}{88.7} % % of circuits where QALM > guoq_2q on fidelity (metrics.win_rate win)
\newcommand{\numWinCxGuoqTwoQ}{55.2} % % of circuits where QALM >= guoq_2q on cx (metrics.win_rate win_or_tie; fidelity ties |d ln F| <= 1e-12) [paper: 54.4 at figure/win_rate.png]
\newcommand{\numStrictWinCxGuoqTwoQ}{46.8} % % of circuits where QALM > guoq_2q on cx (metrics.win_rate win)
\newcommand{\numWinTotalGuoqTwoQ}{95.2} % % of circuits where QALM >= guoq_2q on total (metrics.win_rate win_or_tie; fidelity ties |d ln F| <= 1e-12) [paper: 88.3 at figure/win_rate.png]
\newcommand{\numStrictWinTotalGuoqTwoQ}{92.3} % % of circuits where QALM > guoq_2q on total (metrics.win_rate win)
\newcommand{\numWinFidGuoqTotal}{86.7} % % of circuits where QALM >= guoq_total on fidelity (metrics.win_rate win_or_tie; fidelity ties |d ln F| <= 1e-12) [paper: 95.2 at figure/win_rate.png]
\newcommand{\numStrictWinFidGuoqTotal}{75.8} % % of circuits where QALM > guoq_total on fidelity (metrics.win_rate win)
\newcommand{\numWinCxGuoqTotal}{81.0} % % of circuits where QALM >= guoq_total on cx (metrics.win_rate win_or_tie; fidelity ties |d ln F| <= 1e-12) [paper: 95.6 at figure/win_rate.png]
\newcommand{\numStrictWinCxGuoqTotal}{58.1} % % of circuits where QALM > guoq_total on cx (metrics.win_rate win)
\newcommand{\numWinTotalGuoqTotal}{95.2} % % of circuits where QALM >= guoq_total on total (metrics.win_rate win_or_tie; fidelity ties |d ln F| <= 1e-12) [paper: 96.0 at figure/win_rate.png]
\newcommand{\numStrictWinTotalGuoqTotal}{83.5} % % of circuits where QALM > guoq_total on total (metrics.win_rate win)
\newcommand{\numWinFidQuesoTwoQ}{99.2} % % of circuits where QALM >= queso_2q on fidelity (metrics.win_rate win_or_tie; fidelity ties |d ln F| <= 1e-12) [paper: 98.8 at figure/win_rate.png]
\newcommand{\numStrictWinFidQuesoTwoQ}{94.4} % % of circuits where QALM > queso_2q on fidelity (metrics.win_rate win)
\newcommand{\numWinCxQuesoTwoQ}{96.8} % % of circuits where QALM >= queso_2q on cx (metrics.win_rate win_or_tie; fidelity ties |d ln F| <= 1e-12) [paper: 99.6 at figure/win_rate.png]
\newcommand{\numStrictWinCxQuesoTwoQ}{80.2} % % of circuits where QALM > queso_2q on cx (metrics.win_rate win)
\newcommand{\numWinTotalQuesoTwoQ}{99.6} % % of circuits where QALM >= queso_2q on total (metrics.win_rate win_or_tie; fidelity ties |d ln F| <= 1e-12) [paper: 98.8 at figure/win_rate.png]
\newcommand{\numStrictWinTotalQuesoTwoQ}{94.8} % % of circuits where QALM > queso_2q on total (metrics.win_rate win)
\newcommand{\numWinFidQuesoTotal}{98.8} % % of circuits where QALM >= queso_total on fidelity (metrics.win_rate win_or_tie; fidelity ties |d ln F| <= 1e-12) [paper: 98.4 at figure/win_rate.png]
\newcommand{\numStrictWinFidQuesoTotal}{92.7} % % of circuits where QALM > queso_total on fidelity (metrics.win_rate win)
\newcommand{\numWinCxQuesoTotal}{96.4} % % of circuits where QALM >= queso_total on cx (metrics.win_rate win_or_tie; fidelity ties |d ln F| <= 1e-12) [paper: 99.6 at figure/win_rate.png]
\newcommand{\numStrictWinCxQuesoTotal}{80.6} % % of circuits where QALM > queso_total on cx (metrics.win_rate win)
\newcommand{\numWinTotalQuesoTotal}{99.6} % % of circuits where QALM >= queso_total on total (metrics.win_rate win_or_tie; fidelity ties |d ln F| <= 1e-12) [paper: 98.0 at figure/win_rate.png]
\newcommand{\numStrictWinTotalQuesoTotal}{93.5} % % of circuits where QALM > queso_total on total (metrics.win_rate win)
\newcommand{\numWinFidVoqc}{100.0} % % of circuits where QALM >= voqc on fidelity (metrics.win_rate win_or_tie; fidelity ties |d ln F| <= 1e-12) [paper: 86.5 at figure/win_rate.png]
\newcommand{\numStrictWinFidVoqc}{95.6} % % of circuits where QALM > voqc on fidelity (metrics.win_rate win)
\newcommand{\numWinCxVoqc}{97.6} % % of circuits where QALM >= voqc on cx (metrics.win_rate win_or_tie; fidelity ties |d ln F| <= 1e-12) [paper: 100.0 at figure/win_rate.png]
\newcommand{\numStrictWinCxVoqc}{85.9} % % of circuits where QALM > voqc on cx (metrics.win_rate win)
\newcommand{\numWinTotalVoqc}{100.0} % % of circuits where QALM >= voqc on total (metrics.win_rate win_or_tie; fidelity ties |d ln F| <= 1e-12) [paper: 82.0 at figure/win_rate.png]
\newcommand{\numStrictWinTotalVoqc}{95.6} % % of circuits where QALM > voqc on total (metrics.win_rate win)
\newcommand{\numWinFidQuartz}{99.2} % % of circuits where QALM >= quartz on fidelity (metrics.win_rate win_or_tie; fidelity ties |d ln F| <= 1e-12) [paper: 97.6 at figure/win_rate.png]
\newcommand{\numStrictWinFidQuartz}{91.9} % % of circuits where QALM > quartz on fidelity (metrics.win_rate win)
\newcommand{\numWinCxQuartz}{96.8} % % of circuits where QALM >= quartz on cx (metrics.win_rate win_or_tie; fidelity ties |d ln F| <= 1e-12) [paper: 99.2 at figure/win_rate.png]
\newcommand{\numStrictWinCxQuartz}{78.6} % % of circuits where QALM > quartz on cx (metrics.win_rate win)
\newcommand{\numWinTotalQuartz}{99.6} % % of circuits where QALM >= quartz on total (metrics.win_rate win_or_tie; fidelity ties |d ln F| <= 1e-12) [paper: 96.8 at figure/win_rate.png]
\newcommand{\numStrictWinTotalQuartz}{92.3} % % of circuits where QALM > quartz on total (metrics.win_rate win)
\newcommand{\numWinFidStrongest}{86.7} % lowest fidelity win-or-tie % over the six baselines (the baseline is \numStrongestFidBaseline) [paper: 83.9 at texts/0-abstract.tex:40; texts/1-introduction.tex:136]
\newcommand{\numStrongestFidBaseline}{\guoq{} (Total Gate)} % baseline against which QALM has the lowest fidelity win-or-tie rate [paper: GUOQ (2Q) at texts/0-abstract.tex:39; figtex/gate_scatter.tex:10-11]
\newcommand{\numWinMinOthers}{81.0} % lowest win-or-tie % over QUESO (both), VOQC, Quartz, GUOQ (Total) x {fidelity, cx, total} [paper: 'nearly every circuit' (figure: lowest 82.0, VOQC total) at texts/1-introduction.tex:133-134; figure/win_rate.png]
\newcommand{\numStrictWinFidMin}{75.8} % lowest strict fidelity win % over the six baselines [paper: >50 at texts/5-eval.tex:107-109 ('on the majority of circuits')]
\newcommand{\numNegTotalGuoqTwoQ}{111} % circuits where GUOQ (2Q Gate) output has more gates than the input [paper: many at texts/5-eval.tex:119-120 ('on many circuits')]
\newcommand{\numNegTotalPctGuoqTwoQ}{44.8} % 100 * numNegTotalGuoqTwoQ / n
\newcommand{\numNegTotalQalm}{0} % circuits where the QALM output has more gates than the input [paper: implied 0 ('avoids this entirely') at texts/5-eval.tex:120]

% --- S26 table (tab:eval_end_to_end)
\newcommand{\numGeoRedVoqc}{27.0} % 100*(1 - geomean(opt/orig)) over the 26 circuits, column voqc (S26 originals; metrics.geomean_reduction_pct) [paper: 27.0 at figtex/eval_end_to_end.tex:37]
\newcommand{\numFailETEVoqc}{0} % runs of column voqc whose status is not OK ({})
\newcommand{\numGeoRedQueso}{24.6} % 100*(1 - geomean(opt/orig)) over the 26 circuits, column queso (S26 originals; metrics.geomean_reduction_pct) [paper: 23.9 at figtex/eval_end_to_end.tex:37]
\newcommand{\numFailETEQueso}{0} % runs of column queso whose status is not OK ({})
\newcommand{\numGeoRedGuoq}{28.7} % 100*(1 - geomean(opt/orig)) over the 26 circuits, column guoq (S26 originals; metrics.geomean_reduction_pct) [paper: 29.2 at figtex/eval_end_to_end.tex:37]
\newcommand{\numFailETEGuoq}{0} % runs of column guoq whose status is not OK ({})
\newcommand{\numGeoRedQuartzNoPP}{16.7} % 100*(1 - geomean(opt/orig)) over the 26 circuits, column quartz_nopp (S26 originals; metrics.geomean_reduction_pct) [paper: 16.5 at figtex/eval_end_to_end.tex:37]
\newcommand{\numFailETEQuartzNoPP}{0} % runs of column quartz_nopp whose status is not OK ({})
\newcommand{\numGeoRedQuartzPP}{29.7} % 100*(1 - geomean(opt/orig)) over the 26 circuits, column quartz_pp (S26 originals; metrics.geomean_reduction_pct) [paper: 29.7 at figtex/eval_end_to_end.tex:37]
\newcommand{\numFailETEQuartzPP}{0} % runs of column quartz_pp whose status is not OK ({})
\newcommand{\numGeoRedQalm}{33.9} % 100*(1 - geomean(opt/orig)) over the 26 circuits, column qalm (S26 originals; metrics.geomean_reduction_pct) [paper: 34.3 at figtex/eval_end_to_end.tex:37]
\newcommand{\numFailETEQalm}{0} % runs of column qalm whose status is not OK ({})
\newcommand{\numGeoDeltaVoqc}{6.9} % geomean_red(qalm) - geomean_red(voqc), unrounded, then rounded [paper: 7.3 at texts/5-eval.tex:160]
\newcommand{\numGeoDeltaQuartzNoPP}{17.1} % geomean_red(qalm) - geomean_red(quartz_nopp), unrounded, then rounded [paper: 17.8 at texts/5-eval.tex:162]
\newcommand{\numGeoDeltaQueso}{9.3} % geomean_red(qalm) - geomean_red(queso), unrounded, then rounded [paper: 10.4 at texts/5-eval.tex:163]
\newcommand{\numGeoDeltaGuoq}{5.2} % geomean_red(qalm) - geomean_red(guoq), unrounded, then rounded [paper: 5.1 at texts/5-eval.tex:163]
\newcommand{\numGeoDeltaQuartzPP}{4.1} % geomean_red(qalm) - geomean_red(quartz_pp), unrounded, then rounded [paper: 4.6 at texts/5-eval.tex:164]
\newcommand{\numETEQalmBest}{22} % circuits where QALM has the lowest gate count (ties count) [paper: implied 25 ('all of them except for one') at texts/5-eval.tex:156-158]
\newcommand{\numETEQalmNotBest}{4} % circuits where QALM is not (co-)lowest (26 - numETEQalmBest) [paper: implied 1 ('except for one circuit') at texts/5-eval.tex:158]
\newcommand{\numETEQalmNotBestWord}{four} % numETEQalmNotBest as a word [paper: one at texts/5-eval.tex:158]
\newcommand{\numETEQalmGeqVoqc}{26} % circuits where QALM's gate count <= voqc's
\newcommand{\numETEQalmWorseVoqc}{0} % circuits where QALM's gate count > voqc's
\newcommand{\numETEQalmGeqQueso}{26} % circuits where QALM's gate count <= queso's
\newcommand{\numETEQalmWorseQueso}{0} % circuits where QALM's gate count > queso's
\newcommand{\numETEQalmGeqGuoq}{22} % circuits where QALM's gate count <= guoq's
\newcommand{\numETEQalmWorseGuoq}{4} % circuits where QALM's gate count > guoq's
\newcommand{\numETEQalmGeqQuartzNoPP}{26} % circuits where QALM's gate count <= quartz_nopp's
\newcommand{\numETEQalmWorseQuartzNoPP}{0} % circuits where QALM's gate count > quartz_nopp's
\newcommand{\numETEQalmGeqQuartzPP}{26} % circuits where QALM's gate count <= quartz_pp's [paper: implied 25 ('all but one circuit') at texts/5-eval.tex:171-172]
\newcommand{\numETEQalmWorseQuartzPP}{0} % circuits where QALM's gate count > quartz_pp's [paper: implied 1 ('all but one circuit') at texts/5-eval.tex:171-172]
\newcommand{\numETEQalmWorseQuartzPPWord}{zero} % numETEQalmWorseQuartzPP as a word [paper: one at texts/5-eval.tex:171-172 ('on all but one circuit')]
\newcommand{\numQuartzSuiteCircuits}{26} % circuits in repro/inputs/S26.csv [paper: 26 at texts/5-eval.tex:150,156,331,358,409; figtex/eval_end_to_end.tex:3; ablation figure titles]

% --- greedy mode (fig:greedy_k)
\newcommand{\numGreedyThreshold}{32} % --threshold (mean total reduction %, for time-to-X) [paper: 32 at texts/5-eval.tex:324]
\newcommand{\numGreedyFinalGkZero}{32.25} % final mean total reduction % (t = budget), greedy_k=0, 26 circuits [paper: 32.80 at texts/5-eval.tex:321]
\newcommand{\numGreedyFinalGkOne}{33.12} % final mean total reduction % (t = budget), greedy_k=1, 26 circuits
\newcommand{\numGreedyFinalGkTwo}{32.91} % final mean total reduction % (t = budget), greedy_k=2, 26 circuits [paper: 33.27 at texts/5-eval.tex:322,335]
\newcommand{\numGreedyFinalGkThree}{33.02} % final mean total reduction % (t = budget), greedy_k=3, 26 circuits [paper: 33.12 at texts/5-eval.tex:334]
\newcommand{\numGreedyTimeToThresholdGkZero}{1489} % first t where the mean reduction step function >= 32% (metrics.time_to_reach), greedy_k=0; seconds [paper: 968 at texts/5-eval.tex:323-324]
\newcommand{\numGreedyTimeToThresholdGkOne}{114} % first t where the mean reduction step function >= 32% (metrics.time_to_reach), greedy_k=1; seconds [paper: 89 at texts/5-eval.tex:326]
\newcommand{\numGreedyTimeToThresholdGkTwo}{197} % first t where the mean reduction step function >= 32% (metrics.time_to_reach), greedy_k=2; seconds
\newcommand{\numGreedyTimeToThresholdGkThree}{149} % first t where the mean reduction step function >= 32% (metrics.time_to_reach), greedy_k=3; seconds
\newcommand{\numGreedySpeedupGkOne}{13} % time_to_threshold(greedy_k=0) / time_to_threshold(greedy_k=1) [paper: 11 at texts/5-eval.tex:326]
\newcommand{\numGreedySpeedupGkTwo}{7.6} % time_to_threshold(greedy_k=0) / time_to_threshold(greedy_k=2)
\newcommand{\numGreedySpeedupGkThree}{10} % time_to_threshold(greedy_k=0) / time_to_threshold(greedy_k=3)
\newcommand{\numGreedyFinalDiff}{0.66} % final(greedy_k=2) - final(greedy_k=0), points (unrounded) [paper: <0.5 at texts/5-eval.tex:320-321 ('by less than 0.5\%')]
\newcommand{\numGreedyPhaseMeanGkOne}{7.3} % mean over circuits of the greedy-phase end (first k=2 event; runs that never left it count as the budget), seconds [paper: 5.5 at texts/5-eval.tex:329]
\newcommand{\numGreedyNotFinishedGkOne}{0} % circuits whose greedy phase (greedy_k=1) did not end within the budget
\newcommand{\numGreedyPhaseMeanGkTwo}{107} % mean over circuits of the greedy-phase end (first k=3 event; runs that never left it count as the budget), seconds [paper: 82 at texts/5-eval.tex:329-330]
\newcommand{\numGreedyNotFinishedGkTwo}{0} % circuits whose greedy phase (greedy_k=2) did not end within the budget
\newcommand{\numGreedyPhaseMeanGkThree}{840} % mean over circuits of the greedy-phase end (first k=4 event; runs that never left it count as the budget), seconds [paper: 724 at texts/5-eval.tex:331]
\newcommand{\numGreedyNotFinishedGkThree}{3} % circuits whose greedy phase (greedy_k=3) did not end within the budget [paper: two at texts/5-eval.tex:331]
\newcommand{\numGreedyNotFinishedGkThreeWord}{three} % numGreedyNotFinishedGkThree as a word [paper: two at texts/5-eval.tex:331]
\newcommand{\numGreedyNotFinishedGkThreeList}{\texttt{gf2\^{}8\_mult}, \texttt{gf2\^{}9\_mult}, and \texttt{gf2\^{}10\_mult}} % the circuits counted in numGreedyNotFinishedGkThree [paper: gf2^9_mult and gf2^10_mult at texts/5-eval.tex:332]

% --- fixed exploration depth (fig:single_k)
\newcommand{\numFixedKGkZeroKOne}{31.49} % final mean total reduction %, greedy_k=0, fixed k=1 [paper: 32.83 at texts/5-eval.tex:367]
\newcommand{\numFixedKGkZeroKTwo}{32.36} % final mean total reduction %, greedy_k=0, fixed k=2 [paper: 32.652 (figure only) at figure/single_k_gk0.pdf]
\newcommand{\numFixedKGkZeroKThree}{31.37} % final mean total reduction %, greedy_k=0, fixed k=3 [paper: 32.528 (figure only) at figure/single_k_gk0.pdf]
\newcommand{\numFixedKGkZeroKFour}{31.38} % final mean total reduction %, greedy_k=0, fixed k=4 [paper: 32.100 (figure only) at figure/single_k_gk0.pdf]
\newcommand{\numFixedKGkZeroKFive}{31.24} % final mean total reduction %, greedy_k=0, fixed k=5 [paper: 31.46 at texts/5-eval.tex:368]
\newcommand{\numFixedKGkZeroAdvancing}{32.25} % final mean total reduction %, greedy_k=0, advancing k (same runs as the greedy_k figure)
\newcommand{\numFixedKGkZeroBestK}{2} % fixed k with the highest final reduction, greedy_k=0
\newcommand{\numFixedKGkZeroBest}{32.36} % final reduction of the best fixed k, greedy_k=0
\newcommand{\numFixedKGkZeroGapToAdvancing}{\ensuremath{-}0.12} % advancing-k final - best fixed-k final, greedy_k=0 (points)
\newcommand{\numFixedKGkTwoKOne}{31.60} % final mean total reduction %, greedy_k=2, fixed k=1
\newcommand{\numFixedKGkTwoKTwo}{33.17} % final mean total reduction %, greedy_k=2, fixed k=2
\newcommand{\numFixedKGkTwoKThree}{32.91} % final mean total reduction %, greedy_k=2, fixed k=3
\newcommand{\numFixedKGkTwoKFour}{33.06} % final mean total reduction %, greedy_k=2, fixed k=4
\newcommand{\numFixedKGkTwoKFive}{32.78} % final mean total reduction %, greedy_k=2, fixed k=5
\newcommand{\numFixedKGkTwoAdvancing}{32.91} % final mean total reduction %, greedy_k=2, advancing k (same runs as the greedy_k figure) [paper: 33.27 at texts/5-eval.tex:378]
\newcommand{\numFixedKGkTwoBestK}{2} % fixed k with the highest final reduction, greedy_k=2 [paper: 3 at texts/5-eval.tex:374-375]
\newcommand{\numFixedKGkTwoBest}{33.17} % final reduction of the best fixed k, greedy_k=2 [paper: 33.11 at texts/5-eval.tex:375]
\newcommand{\numFixedKGkTwoGapToAdvancing}{\ensuremath{-}0.26} % advancing-k final - best fixed-k final, greedy_k=2 (points)

% --- exploration strategy (fig:exploration_ablation)
\newcommand{\numNpoolOne}{32.94} % final mean total reduction %, N_pool=1, mean over seeds [0, 1, 2] [paper: 33.16 at figure/N_pool_arith_mean_3600s_plain.pdf; texts/5-eval.tex:419 (range end)]
\newcommand{\numNpoolTwo}{32.65} % final mean total reduction %, N_pool=2, mean over seeds [0, 1, 2]
\newcommand{\numNpoolThree}{32.68} % final mean total reduction %, N_pool=3, mean over seeds [0, 1, 2]
\newcommand{\numNpoolFour}{32.68} % final mean total reduction %, N_pool=4, mean over seeds [0, 1, 2] [paper: 33.03 at figure/N_pool_arith_mean_3600s_plain.pdf; texts/5-eval.tex:419 (range end)]
\newcommand{\numNpoolFive}{32.69} % final mean total reduction %, N_pool=5, mean over seeds [0, 1, 2]
\newcommand{\numNbranchOne}{31.51} % final mean total reduction %, N_branch=1, mean over seeds [0, 1, 2] [paper: 31.33 at texts/5-eval.tex:431]
\newcommand{\numNbranchTwo}{32.36} % final mean total reduction %, N_branch=2, mean over seeds [0, 1, 2] [paper: 32.75 at texts/5-eval.tex:433]
\newcommand{\numNbranchThree}{32.94} % final mean total reduction %, N_branch=3, mean over seeds [0, 1, 2] [paper: 33.37 at texts/5-eval.tex:434]
\newcommand{\numNbranchFour}{32.91} % final mean total reduction %, N_branch=4, mean over seeds [0, 1, 2] [paper: 33.30 at texts/5-eval.tex:437]
\newcommand{\numNbranchFive}{33.11} % final mean total reduction %, N_branch=5, mean over seeds [0, 1, 2] [paper: 33.21 at texts/5-eval.tex:438]
\newcommand{\numNpoolMin}{32.65} % mean_final_min_pct over N_pool in 1..5 (seed-averaged finals) [paper: 33.03 at texts/5-eval.tex:418-419]
\newcommand{\numNpoolMax}{32.94} % mean_final_max_pct over N_pool in 1..5 (seed-averaged finals) [paper: 33.16 at texts/5-eval.tex:418-419]
\newcommand{\numNpoolSpread}{0.29} % mean_final_spread_pct over N_pool in 1..5 (seed-averaged finals) [paper: 0.13 at texts/5-eval.tex:418-419]
\newcommand{\numNbranchSpread}{1.59} % max - min over N_branch in 1..5 (seed-averaged finals)
\newcommand{\numSeeds}{3} % seeds per E5 configuration (fig_curves.E5_SEEDS)
\newcommand{\numSeedSdDefault}{0.12} % sample SD over seeds of the final mean reduction, default config (N_pool=1, N_branch=3)
\newcommand{\numSeedRangeDefault}{0.24} % max - min over seeds of the final mean reduction, default config
\newcommand{\numSeedSdPooled}{0.09} % sqrt(pooled within-config variance over seeds) of the final mean reduction, over the distinct E5 configurations
\newcommand{\numSeedHsd}{0.22} % Tukey HSD (95%) for 5 means of len(seeds) runs: q(0.95; 5, df) * SD_pooled / sqrt(seeds)
